\section{Local prediction: protocols, controls, and robustness}
\label{app:local_evidence}

This appendix follows the two local studies in Section~\ref{sec:exp_local}:
archived Pinball prediction first, then Circular architecture controls.
Their model identities, windows, and sampling protocols are distinct;
their numerical errors are not pooled.

\subsection{Pinball: archived local-prediction protocol}
\label{app:pinball}

\paragraph{Predictor provenance.}
The Periodic rows in Table~\ref{tab:pinball_autonomous} are associated
with Deep-FNN-H3. The formal outer-candidate records also identify the
three Pinball assembly candidates as Deep-FNN-H3, but these records do
not establish the architecture of every earlier standalone Steady/Hopf
row. The table therefore uses a generic archived-predictor label for
those rows; no numerical values have been reassigned to a MoE model.

\paragraph{Periodic-core evaluation protocol.}
The $K=32$ Periodic result in Table~\ref{tab:pinball_autonomous} corresponds to the standard periodic evaluation. We additionally report a longer $K=56$ evaluation restricted to the Periodic-core final-test trajectories. The full-order snapshots are stored every $0.25$ physical time units; subsampling every fourth snapshot gives an effective interval $\Delta t=1$, matching the Periodic specialist's native history cadence. The $K=56$ setting therefore provides a longer autonomous evaluation within the same periodic regime without changing the specialist representation or
sampling convention.

\paragraph{Separate one-step versus autonomous diagnostic.}
The POD--DeepONet baseline provides a complementary comparison between
one-step prediction and autonomous rollout. It attains $E_u=1.5921\%$ and $E_p=23.3447\%$ on 4,228 held-out input--output pairs, but only 1 of 28 held-out trajectories completes the prescribed autonomous rollout evaluation. This difference illustrates that one-step accuracy does not by itself imply reliable autonomous prediction. Because the POD--DeepONet evaluation uses a different temporal cadence from
the Periodic specialist, it is used only for this comparison of one-step prediction and autonomous rollout and not as a cadence-matched long-horizon accuracy baseline.

\subsection{Circular: controlled local-model protocol}
\label{app:circular}

\paragraph{Evaluation and aggregation protocol.}
The neural study uses seeds 1248, 1600, and 2026.
Hopf has 42 windows from
$Re=47.081355,49.022357,51.786450$; Periodic has 32 windows from
$Re=70.314635,100.352251,149.059229,189.862278$.
Each window contains 48 prediction steps, with $K=24$ using its
first 24 steps. Relative physical-field errors are averaged over
steps, then windows within each Reynolds number, then equally over
Reynolds numbers. Means and sample standard deviations are computed
over accepted seeds. Reference fields are reconstructed from true
POD coefficients on the corresponding basis.
Only checkpoints accepted by the original validation rule are
evaluated on held-out windows.
Table~\ref{tab:specialist_training} describes the archived original
training configurations, not the selected checkpoints of every
repeated run.

Dense matches active neural parameter count, not total stored
parameters or execution time. Structured removes input-dependent
expert routing while retaining explicit response branches.
Neither is identical to the older Vanilla-FNN-MoE or Global MoE
configuration. The present protocol does not retroactively validate
the earlier Circular min--max summaries or establish the identity
of the unresolved Steady comparison.
Joint errors and their interpretation are reported once in
Table~\ref{tab:circular_full_reorg} and the accompanying main-text
discussion; the supplementary table below supplies component errors.

\begin{table}[t]
\centering
\small
\setlength{\tabcolsep}{3.5pt}
\caption{Circular Periodic component errors (\%). Neural entries
are mean $\pm$ sample standard deviation over three accepted seeds.
OpInf is deterministic. Joint errors are reported in the main-text
comparison and are not repeated here. Components and joint values
are independently rounded.}
\label{tab:circular_full}
\begin{tabular}{@{}lcccc@{}}
\toprule
& \multicolumn{2}{c}{$K=24$} & \multicolumn{2}{c}{$K=48$}\\
\cmidrule(lr){2-3}\cmidrule(l){4-5}
Model & $E_u$ & $E_p$ & $E_u$ & $E_p$\\
\midrule
Proposed, original
& $0.5731\pm0.3638$ & $2.3060\pm1.3681$
& $1.0107\pm0.6513$ & $3.9337\pm2.5070$\\
Proposed, repaired
& $0.5814\pm0.3949$ & $2.4854\pm1.4427$
& $1.1040\pm0.7840$ & $4.4282\pm2.8786$\\
Dense, active-matched
& $0.2446\pm0.0323$ & $1.0805\pm0.1946$
& $0.4790\pm0.1208$ & $1.9208\pm0.5391$\\
Structured, original
& $0.5178\pm0.2667$ & $2.3489\pm1.1688$
& $1.0405\pm0.5015$ & $4.2406\pm2.0664$\\
Structured, repaired
& $0.5429\pm0.2783$ & $2.3168\pm1.0315$
& $0.9983\pm0.4944$ & $4.0693\pm1.8827$\\
OpInf & 15.8676 & 98.2264 & 17.6639 & 97.3258\\
\bottomrule
\end{tabular}
\end{table}

\paragraph{Quadratic-factor initialization.}
For a multiplicative response
$(\bm v_1^\top\bm s)(\bm v_2^\top\bm s)$, the gradient of either
factor contains the other factor. Initializing both at zero therefore
prevents learning in this branch.
Each accepted original Hopf Proposed checkpoint contains 28
quadratic-factor tensors, all zero. Each original Periodic Proposed
checkpoint contains 84, and each Periodic Structured checkpoint
contains four, also all zero.
The repaired initialization preserves a random factor and zeros
only its partner. Smoke training updates all 14 factor pairs for
Hopf Proposed, 42 for Periodic Proposed, and two for Periodic
Structured. This check measures branch learning activity, separately
from the accuracy comparison in the main text.

\paragraph{Training completion and validation acceptance.}
Table~\ref{tab:hopf_completion} records the attempted Hopf campaigns.
Acceptance requires a usable run and the prescribed validation rule;
a missing accepted checkpoint is not automatically a divergent
held-out rollout. The original Dense control also has two accepted
seeds out of three.

\begin{table}[t]
\centering
\small
\caption{Hopf completion and validation acceptance over three
attempted seeds per architecture. AMP and FP32 are separate repaired
campaigns, not pooled repeats.}
\label{tab:hopf_completion}
\begin{tabular}{@{}lccp{5.8cm}@{}}
\toprule
Configuration & Proposed & Structured & Failure or rejection\\
\midrule
Original & 2/3 & 3/3 & Proposed seed 1600 not accepted by validation\\
Repaired, AMP & 0/3 & 0/3 & Non-finite gradients\\
Repaired, FP32 & 2/3 & 2/3 & Proposed 2026: step 3220;
Structured 1248: step 2860, both during $K=4$\\
\bottomrule
\end{tabular}
\end{table}

The FP32 preflight covers 116 combinations of 29 training parameters
and four curriculum lengths. Passing it does not guarantee finite
gradients later in optimization.
The repaired FP32 training-failure rate is $1/3$ per architecture
($2/6$ pooled across these two configurations). The earlier AMP
campaign failed for all six attempts and is not pooled with FP32.
Original validation rejection is distinct from non-finite-gradient
training failure. Failed seeds are not replaced by post-hoc retries
in the reported statistics.

\paragraph{OpInf adaptation and validation.}
The regularized polynomial model uses constant, linear, compact
quadratic, standardized $1/Re$, and parameter--state interaction
features. Training-trajectory derivatives are estimated by
second-order numerical differentiation.
Ridge regularization is selected on validation windows, with a
non-finite window invalidating a candidate.
Prediction uses fixed RK4 with four substeps per output interval.
Pressure is an algebraic ridge output using the same polynomial
features rather than the proposed PPE-plus-correction map.
Test windows, reference reconstructions, and quadrature match the
neural comparison. The upper endpoint of the Periodic ridge search
is 100. This is a task-specific adaptation, not a reproduction of
every implementation detail of an external method.
